\section{The Question Before the Preference}

Preference is often declared before the liturgical and theological question is faced. A comparison that rewards or punishes a Missal for practices not in it cannot reach a just conclusion. The first task is therefore to ask what sacrifice each received book enacts, how each orders participation, and which differences belong to the book, governing law, or local practice.

\newcommand{\missalscopeappendix}{%
\section{Scope, Edition, and Qualifications}

\subsection{Exact comparison witnesses}

The controversy is usually named before its objects are defined. “The TLM” may mean the Roman Missal of 1962, an earlier Holy Week, the pre-1955 calendar, any Latin Mass before 1970, or an entire social and devotional world. “The Novus Ordo” may mean Paul VI's 1969 Order of Mass, the complete Latin Missal of 1970, the third typical edition, an English translation, celebration facing the people, Communion in the hand, guitars, altar girls, or an abuse encountered in one parish. Those bundles must be disaggregated before comparison.

This paper compares two exact normative witnesses: the 1962 Latin typical edition of the \emph{Missale Romanum} promulgated under John~XXIII, and the postconciliar Roman Missal represented by its 2002 third Latin typical edition and 2008 emended reprint. Where concrete territorial discipline requires it, the current United States English Missal, General Instruction, and complementary norms control. Latin controls the identity of both typical editions; an approved U.S. English text controls only its identified territorial use. Earlier, transitional, and other Latin editions or uses are identified in their own right rather than blended into either witness.

“Traditional Latin Mass” is retained as a familiar name for celebration according to the 1962 Missal, not as a claim that it exhausts Latin tradition. “Postconciliar Missal” is normally preferred to “Novus Ordo,” because the latter can ambiguously name the 1969 Order, the 1970 book, a later typical edition, a vernacular translation, or local practice. The historical comparison follows antecedent Roman strata and printed witnesses through the 1962 book, the conciliar reform and transitional books, and the 2002/2008 postconciliar horizon. Historical witness checks extend through 17 July 2026; mutable law, permissions, and the community census are frozen at 13 July 2026 unless a claim states an earlier date.

\subsection{Five questions that must not collapse}

\begin{longtable}{>{\bfseries\raggedright\arraybackslash}p{0.16\linewidth}>{\raggedright\arraybackslash}p{0.34\linewidth}>{\raggedright\arraybackslash}p{0.42\linewidth}}
\toprule
\textbf{Question} & \textbf{What answers it} & \textbf{What does not answer it}\\
\midrule
\endhead
Validity & Matter, sacramental form, valid minister, ecclesial intention, and a capable subject where the sacrament requires one. & Beauty, antiquity, vernacular language, a minister's personal holiness, or a congregation's fervor.\\
Liceity & Universal and particular law, faculties, delegation, office, permissions, and faithful observance of the approved rite. & Validity alone, private necessity theories, or the claim that a preferred form is holier.\\
Doctrinal integrity & The official rite read as a whole in the Church's faith and magisterial interpretation. & An abuse, a tendentious paraphrase, or the absence of every possible doctrinal emphasis from every prayer.\\
Historical continuity & Textual ancestry, ritual structure, development, reception, authority, and the scale and manner of change. & The bare fact of validity or a slogan that nothing changed because the same sacrifice remains.\\
Pastoral prudence & Reverence, intelligibility, formation, stability, local circumstances, genuine participation, and observed fruits under Catholic doctrine. & A universal ranking deduced from one parish, one generation, or one temperament.\\
\bottomrule
\end{longtable}

Validity is the floor beneath Catholic comparison, not its ceiling. If an enacted Mass truly makes present Christ's one sacrifice, it is not thereby shown that every textual choice, translation, option, architectural fashion, or manner of distribution is equally wise. Conversely, a severe historical or pastoral criticism cannot be converted into sacramental invalidity merely to make the criticism sound ultimate. The Church's categories protect both truth and the freedom to reason.

\missalevidenceappendix
}

\subsection{The governing judgment}

The Roman Missal of 1962 is not “the Mass invented at Trent.” It is the last preconciliar typical edition of a Roman curial and Franciscan inheritance substantially visible before Trent, standardized by Pius V in 1570, and repeatedly revised by his successors. The 1570 law displaced younger local variants of the Roman custom while expressly protecting qualifying ancient uses; it neither abolished every non-Roman Latin liturgy nor furnished a name-by-name list of suppressed books.

The postconciliar Missal is not “the Mass written by Vatican II.” The council mandated a general reform under limiting principles. A papal commission then produced, and Paul VI personally authorized and promulgated, a source-rich twentieth-century reconstruction and expansion of the Roman books. It retains the sacrificial and sacramental identity of the Roman Mass while materially changing the lectionary, calendar, presidential prayers, offertory, Eucharistic-prayer repertoire, ceremonial, ministries, and experience of variation.

Both Missals can be celebrated validly, reverently, and fruitfully; both belong to the Catholic history of the Roman Rite. Their equality in sacramental reality does not make their ritual forms interchangeable. The 1962 book carries a denser continuity with the late-medieval and Tridentine inheritance, stronger fixed ascetical and propitiatory expression, and a stable hieratic dramatic form. The postconciliar book offers a much wider scriptural proclamation, more audible common prayer, a larger ancient and newly composed euchological corpus, and differentiated lay participation, but also exposes celebration to more options, presidential visibility, local variation, and discontinuity of received ritual memory. These are comparative findings, not rival creeds.

\governingjudgment{The Catholic answer is not “old valid, new invalid,” nor “new authorized, therefore every criticism disloyal.” One sacrifice stands beneath two historically unequal but ecclesially authorized Missals. Truth requires the Church's sacramental judgment; charity requires refusing caricature; prudence requires naming both the reform's real gains and its real costs.}

\newcommand{\missalevidenceappendix}{%
\subsection{Evidence, corpus, and method}

Official books control what each rite actually prescribes. Conciliar and papal acts control doctrine and promulgation. The Code and competent dicasteries control current universal discipline. Historical reconstruction uses primary witnesses where they survive and specialist scholarship where no decree can narrate development. A community's own site may establish what it does; only competent ecclesiastical authority can establish its juridic status or faculties.

Direct quotations are short. Reproduced liturgical wording comes from identified received witnesses; surrounding English is analytical summary, not recitation text. Claims of “suppression,” “abrogation,” “right,” “indult,” and “schism” retain their legal objects and dates. The community appendix includes uncertain rows rather than laundering absence of evidence into approval. Its audit rule is published because no public universal registry answers the popular request for “all TLM congregations.”

Finally, this is a Catholic comparison, not an external plebiscite on dogma. It affirms Trent's doctrine of sacrifice and Real Presence, the Church's authority over sacramental rites short of their divinely instituted substance, Vatican II's legitimate authority, and the validity of the revised rites. It also applies \emph{Sacrosanctum Concilium}'s own demand that innovation be genuinely required and grow organically from existing forms. The Church's act and the Church's stated criterion belong together.

The corpus covers pre-Tridentine Roman diversity and the printed antecedents of 1570; revision through 1962; conciliar and postconciliar reform; direct structural comparison; sacramental validity and Eucharistic discipline; current older-book law; a bounded corporate-community census; and pastoral judgment. This source-audited working paper is not an act of the Magisterium, a parallel Missal, a ceremonial, a liturgical authorization, an official canonical census, an exhaustive register of Mass sites, or an opinion on an individual's Orders, faculties, culpability, sacramental eligibility, or canonical status. Short received wording identifies an exact witness; the paper neither republishes vernacular liturgical books nor supplies recitation text.

\subsection{Legal scope and currentness}

The current-law analysis applies the 1983 Code of Canon Law and universal Latin-Church liturgical law, with publicly located special provisions, through 13 July 2026. It does not apply the Eastern Code to Eastern Catholics, inventory every diocesan decree, infer a private rescript's terms, or replace advice from the competent authority or a qualified canonist with the complete facts. No act of Leo~XIV repealing or replacing \emph{Traditionis custodes} was located through the cutoff. The controlling general framework is therefore the 2021 motu proprio, its 2021 \emph{Responsa}, and the 2023 rescript, subject to valid particular, proper, or special law. Later competent law supersedes this snapshot.

\subsection{Audit and review limits}

The primary acts, edition boundaries, sacramental categories, and current legal endpoint have been source-audited. The community census is reproducibly bounded but intrinsically perishable. Independent review by a liturgical historian, sacramental theologian, Latinist, canonist, and persons competent in the proper law of the listed institutes remains outstanding.
}
